\documentclass[11pt,a4paper]{scrartcl}

\usepackage{amssymb}
\usepackage{amsmath}

\usepackage[utf8]{inputenc}
\usepackage[T1]{fontenc}
\newcommand{\ats}[1]{\par\noindent\textit{Atsakymas:} #1}
\usepackage{cancel}
\usepackage{tikz}
\usepackage{lmodern} 
\usepackage{amsmath,amssymb,amsthm}
\usepackage{graphicx}
\usepackage[dvipsnames,svgnames]{xcolor}
\usepackage{hyperref}
\usepackage{mathrsfs}
\usepackage[shortlabels]{enumitem}
\usepackage[obeyFinal,textsize=scriptsize,shadow,loadshadowlibrary]{todonotes}
\usepackage{mathtools}
\usepackage{microtype}

%%% Paragraph layout
\setlength{\parindent}{0pt}
\setlength{\parskip}{0.6\baselineskip}

%%% Colored sections
\renewcommand*{\sectionformat}%
  {\color{purple}\S\thesection\enskip}
\renewcommand*{\subsectionformat}%
  {\color{purple}\S\thesubsection\enskip}
\renewcommand*{\subsubsectionformat}%
  {\color{purple}\S\thesubsubsection\enskip}
\addtokomafont{paragraph}{\color{orange!35!black}\P\ }
\KOMAoptions{numbers=noenddot}

%%% Title
\addtokomafont{subtitle}{\Large}
\setkomafont{author}{\Large\scshape}
\setkomafont{date}{\Large\normalsize}
% Neigiamas vertikalus tarpas, kuris pakelia dokumento antraštę į viršų. 
% Galite keisti -2.5cm į -3cm (aukščiau) arba -1.5cm (žemiau).
\titlehead{\vspace*{-7.5cm}} 

% PRIDĖTA: Automatiškai perrašome datą į mokyklos pavadinimą
\AtBeginDocument{%
  \date{\normalsize Vilniaus licėjus}
}

%%% Page setup
\usepackage[headsepline]{scrlayer-scrpage}
\renewcommand{\headfont}{}
\addtolength{\textheight}{3.14cm}
\setlength{\footskip}{0.5in}
\setlength{\headsep}{10pt}
% Puslapio antraštė turi rodyti tik konkrečius dokumento duomenis.
% Nustatome juos aiškiai, kad neatsirastų "\@author" / "\@title" arba senų reikšmių.
\makeatletter
\AtBeginDocument{%
  \ihead{\footnotesize\textbf{Andrius Berniukevičius}}%
  \ohead{\footnotesize\textbf{\@title}}%
}
\makeatother
\automark{section}
\chead{}
\cfoot{\pagemark}

%%% Macros
\providecommand{\ol}{\overline}
\providecommand{\ul}{\underline}
\providecommand{\wt}{\widetilde}
\providecommand{\wh}{\widehat}
\providecommand{\eps}{\varepsilon}
\providecommand{\half}{\frac{1}{2}}
\providecommand{\inv}{^{-1}}
\newcommand{\dang}{\measuredangle} %% Directed angle
\newcommand{\vocab}[1]{\textbf{\color{blue}\sffamily #1}}
\providecommand{\CC}{\mathbb C}
\providecommand{\FF}{\mathbb F}
\providecommand{\NN}{\mathbb N}
\providecommand{\QQ}{\mathbb Q}
\providecommand{\RR}{\mathbb R}
\providecommand{\ZZ}{\mathbb Z}
\providecommand{\ts}{\textsuperscript}
\providecommand{\dg}{^\circ}
\providecommand{\ii}{\item}
\providecommand{\defeq}{\coloneqq}
\DeclareMathOperator*{\lcm}{lcm}
\DeclareMathOperator*{\argmin}{arg min}
\DeclareMathOperator*{\argmax}{arg max}
\providecommand{\hrulebar}{\par
\hspace{\fill}\rule{0.95\linewidth}{.7pt}\hspace{\fill}
\par\nointerlineskip \vspace{\baselineskip}}

%%% Optional colored theorems
\usepackage{thmtools}
\usepackage[framemethod=TikZ]{mdframed}
\mdfdefinestyle{mdbluebox}{%
  roundcorner=10pt,
  linewidth=1pt,
  skipabove=12pt,
  innerbottommargin=9pt,
  skipbelow=2pt,
  linecolor=blue,
  nobreak=true,
  backgroundcolor=TealBlue!5,
}
\declaretheoremstyle[
  headfont=\sffamily\bfseries\color{MidnightBlue},
  mdframed={style=mdbluebox},
  headpunct={\\[3pt]},
  postheadspace={0pt}
]{thmbluebox}
\mdfdefinestyle{mdredbox}{%
  linewidth=0.5pt,
  skipabove=12pt,
  frametitleaboveskip=5pt,
  frametitlebelowskip=0pt,
  skipbelow=2pt,
  frametitlefont=\bfseries,
  innertopmargin=4pt,
  innerbottommargin=8pt,
  nobreak=true,
  backgroundcolor=Salmon!5,
  linecolor=RawSienna,
}
\declaretheoremstyle[
  headfont=\bfseries\color{RawSienna},
  mdframed={style=mdredbox},
  headpunct={\\[3pt]},
  postheadspace={0pt},
]{thmredbox}
\mdfdefinestyle{mdgreenbox}{%
  skipabove=8pt,
  linewidth=2pt,
  rightline=false,
  leftline=true,
  topline=false,
  bottomline=false,
  linecolor=ForestGreen,
  backgroundcolor=ForestGreen!5,
}
\declaretheoremstyle[
  headfont=\bfseries\sffamily\color{ForestGreen!70!black},
  bodyfont=\normalfont,
  spaceabove=2pt,
  spacebelow=1pt,
  mdframed={style=mdgreenbox},
  headpunct={ --- },
]{thmgreenbox}
\mdfdefinestyle{mdblackbox}{%
  skipabove=8pt,
  linewidth=3pt,
  rightline=false,
  leftline=true,
  topline=false,
  bottomline=false,
  linecolor=black,
  backgroundcolor=RedViolet!5!gray!5,
}
\declaretheoremstyle[
  headfont=\bfseries,
  bodyfont=\normalfont\small,
  spaceabove=0pt,
  spacebelow=0pt,
  mdframed={style=mdblackbox}
]{thmblackbox}

%% Aplinkos su rėmeliais (thmtools)
\declaretheorem[style=thmbluebox,name=Teorema]{theorem}
\declaretheorem[style=thmbluebox,name=Lema]{lemma}
\declaretheorem[style=thmbluebox,name=Savybė]{property}
\declaretheorem[style=thmbluebox,name=Teiginys]{proposition}
\declaretheorem[style=thmbluebox,name=Išvada]{corollary}
\declaretheorem[style=thmbluebox,name=Teorema,numbered=no]{theorem*}
\declaretheorem[style=thmbluebox,name=Lema,numbered=no]{lemma*}
\declaretheorem[style=thmbluebox,name=Teiginys,numbered=no]{proposition*}
\declaretheorem[style=thmbluebox,name=Išvada,numbered=no]{corollary*}
\declaretheorem[style=thmgreenbox,name=Algoritmas]{algorithm}
\declaretheorem[style=thmgreenbox,name=Algoritmas,numbered=no]{algorithm*}
\declaretheorem[style=thmgreenbox,name=Teiginys]{claim}
\declaretheorem[style=thmgreenbox,name=Teiginys,numbered=no]{claim*}
\declaretheorem[style=thmredbox,name=Pavyzdys]{example}
\declaretheorem[style=thmredbox,name=Pavyzdys,numbered=no]{example*}
\declaretheorem[style=thmblackbox,name=Pastaba]{remark}
\declaretheorem[style=thmblackbox,name=Pastaba,numbered=no]{remark*}

%% Standartinės amsthm aplinkos (Apibrėžimai ir užduotys)
\theoremstyle{definition}
\ifcsname conjecture\endcsname\else\newtheorem{conjecture}{Hipotezė}\fi
\ifcsname definition\endcsname\else\newtheorem{definition}{Apibrėžimas}\fi
\ifcsname fact\endcsname\else\newtheorem{fact}{Faktas}\fi
\ifcsname answer\endcsname\else\newtheorem{answer}{Atsakymas}\fi
\ifcsname ques\endcsname\else\newtheorem{ques}{Klausimas}\fi
\ifcsname exercise\endcsname\else\newtheorem{exercise}{Užduotis}\fi
\ifcsname problem\endcsname\else\newtheorem{problem}{Uždavinys}[section]\fi
\ifcsname conjecture*\endcsname\else\newtheorem*{conjecture*}{Hipotezė}\fi
\ifcsname definition*\endcsname\else\newtheorem*{definition*}{Apibrėžimas}\fi
\ifcsname fact*\endcsname\else\newtheorem*{fact*}{Faktas}\fi
\ifcsname answer*\endcsname\else\newtheorem*{answer*}{Atsakymas}\fi
\ifcsname ques*\endcsname\else\newtheorem*{ques*}{Klausimas}\fi
\ifcsname exercise*\endcsname\else\newtheorem*{exercise*}{Užduotis}\fi
\ifcsname problem*\endcsname\else\newtheorem*{problem*}{Uždavinys}\fi

\newcounter{answerssection}
\newcommand{\answerssection}{%
  \setcounter{answerssection}{\value{section}}%
  \section{Atsakymai}%
  \setcounter{problem}{0}%
  \renewcommand{\theproblem}{\arabic{answerssection}.\arabic{problem}}%
}

%% GALUTINIS SPRENDIMAS PROOF APLINKAI
%% --- PRADŽIA: Įrodymo aplinkos sutvarkymas ---
\makeatletter

% 1. Užtikriname, kad žodis būtų "Įrodymas", net jei "babel" paketas bando jį perrašyti
\AtBeginDocument{%
  \renewcommand{\proofname}{Įrodymas}%
  \@ifpackageloaded{babel}{%
    \addto\captionslithuanian{\renewcommand{\proofname}{Įrodymas}}%
  }{}%
}

% 2. Priverstinai pašaliname scrartcl klasės bold šriftą ir paliekame TIK italic
% 3. Sujungiame su tuščiu kvadratėliu dešiniajame kampe.
\renewcommand{\qedsymbol}{\ensuremath{\Box}}
\renewenvironment{proof}[1][\proofname]{\par
  \pushQED{\qed}%
  \normalfont \topsep6\p@\@plus6\p@\relax
  \trivlist
  \item[\hskip\labelsep
        \normalfont\itshape % Priverčia naudoti tik pasvirą šriftą, kaip nuotraukoje
    #1\@addpunct{.}]\ignorespaces
}{%
  \popQED\endtrivlist\@endpefalse
}
\newenvironment{solution}{\par
  \normalfont \topsep6\p@\@plus6\p@\relax
  \trivlist
  \item[\hskip\labelsep
        \normalfont\itshape
    \textit{Sprendimas}\@addpunct{.}]\ignorespaces
}{%
  \endtrivlist\@endpefalse
}
\newcommand{\ans}[1]{\par\noindent\textit{Atsakymas:} #1}
\makeatother


\title{Algebra: skaidymas dauginamaisiais}
\author{Andrius Berniukevičius}

\begin{document}

\maketitle

\section{Skiriamoji riba: atskliausti ar sutraukti?}

Mokyklinėje algebroje jūs nuolat atliekate vieną ir tą patį veiksmą – \textbf{atskliaudžiate}. Pamatę išraišką $(x+2)(x-3)$, jūs iškart ją dauginate ir gaunate $x^2 - x - 6$. Tai yra „griovimo“ procesas. 

Olimpiadinėje matematikoje viskas veikia atvirkščiai. Laimi tas, kuris sugeba \textbf{statyti} struktūras, o ne jas griauti. Mūsų tikslas yra paimti palaidą skaičių ar raidžių kratinį ir paversti jį \textbf{dauginamaisiais} (skliaustais). Kodėl? Nes kai reiškinys tampa sandauga (pavyzdžiui, lygi nuliui, arba nagrinėjant jos dalumą), mes galime daryti griežtas logines išvadas.

\subsection*{Olimpiadininko ginklų arsenalas}
Šias formules jūs žinote mintinai. Tačiau dabar privalote išmokti jas pamatyti net ten, kur jos gudriai paslėptos.

\begin{enumerate}
    \item \textbf{Kvadratų skirtumas (Absoliutus lyderis):} $a^2 - b^2 = (a-b)(a+b)$
    \item \textbf{Pilnas kvadratas (Suma ir skirtumas):} $a^2 \pm 2ab + b^2 = (a \pm b)^2$
    \item \textbf{Kubų suma ir skirtumas:} $a^3 \pm b^3 = (a \pm b)(a^2 \mp ab + b^2)$
\end{enumerate}

% ==========================================
% PRIDĖK IR ATIMK TRIUKAS
% ==========================================
\section{Magija: Triukas „Pridėk ir Atimk“}

Kaip išskaidyti reiškinį $x^4 + x^2 + 1$? 
Jokios standartinės formulės čia tiesiogiai netinka. Tačiau įsižiūrėkime atidžiau. Mes turime $x^4$ (tai yra $(x^2)^2$) ir turime $1$ (tai yra $1^2$). Trūksta tik vidurinio nario, kad gautume pilną kvadratą. Pilnam kvadratui $(x^2 + 1)^2$ per vidurį reikia nario $2x^2$. O mes turime tik $1x^2$.

Matematikoje galioja nuostabi taisyklė: \textbf{prie reiškinio visada galima pridėti nulį}. 
O kas yra nulis? Tai yra $+x^2 - x^2$. 
Pridėkime ir atimkime $x^2$:
\[ x^4 + x^2 + 1 \quad = \quad x^4 + \mathbf{2x^2} + 1 - \mathbf{x^2} \]
Pirmi trys nariai dabar idealiai susiskleidžia į pilną kvadratą!
\[ (x^2 + 1)^2 - x^2 \]
Ką gavome? Tobulą \textbf{kvadratų skirtumą} ($A^2 - B^2$)! Išskaidome:
\[ (x^2 + 1 - x)(x^2 + 1 + x) \]
Perrašome gražiau (pagal laipsnių mažėjimą):
\[ \mathbf{(x^2 - x + 1)(x^2 + x + 1)} \]
Išskaidyta! Šis „Pridėk ir Atimk“ metodas yra vienas galingiausių įrankių olimpiadose.

% ==========================================
% SOPHIE GERMAIN TAPATYBĖ
% ==========================================
\section{Sophie Germain Tapatybė}

XVIII amžiuje gyvenusi prancūzų matematikė Sophie Germain atrado vieną gražiausių ir dažniausiai olimpiadose pasitaikančių tapatybių. Kaip išskaidyti sumą $a^4 + 4b^4$? (Prisiminkite, kad kvadratų \textit{sumai} formulės nėra!)

Pritaikysime tą patį triuką. Turime $(a^2)^2$ ir $(2b^2)^2$. Kad turėtume pilną kvadratą $(a^2 + 2b^2)^2$, mums trūksta vidurinio nario: $2 \cdot a^2 \cdot 2b^2 = 4a^2b^2$. 
Pridėkime jį ir iškart atimkime!
\[ a^4 + 4b^4 \quad = \quad a^4 + \mathbf{4a^2b^2} + 4b^4 - \mathbf{4a^2b^2} \]
Grupuojame pirmus tris narius į pilną kvadratą:
\[ (a^2 + 2b^2)^2 - 4a^2b^2 \]
Pastebime, kad $4a^2b^2$ taip pat yra kvadratas: $(2ab)^2$. Gauname kvadratų skirtumą:
\[ (a^2 + 2b^2)^2 - (2ab)^2 \]
Pritaikome kvadratų skirtumo formulę ir gauname legendinę \textbf{Sophie Germain tapatybę}:
\begin{center}
\fbox{
\begin{minipage}{0.9\textwidth}
\textbf{Sophie Germain tapatybė:}
\[ a^4 + 4b^4 = (a^2 - 2ab + 2b^2)(a^2 + 2ab + 2b^2) \]
\end{minipage}
}
\end{center}

% ==========================================
% PAVYZDŽIAI
% ==========================================
\section{Olimpiadiniai pavyzdžiai}

\begin{example}[Grupavimo menas ir paslėptos formulės]
Išskaidykite dauginamaisiais reiškinius:
a) $x^3 - 3x^2 - 4x + 12$
b) $a^2 - b^2 + 5a - 5b$
\end{example}

\textbf{Sprendimas:}\\
\textbf{a)} Kai matome 4 narius, dažniausiai veikia grupavimas poromis (2 ir 2). Sugrupuokime pirmą narį su antru, o trečią – su ketvirtu. \textbf{Atsargiai su minuso ženklu!}
\[ (x^3 - 3x^2) - (4x - 12) \]
\textit{(Atkreipkite dėmesį: prieš antrą skliaustą iškėlus minusą, skliausto viduje plius $12$ pavirto į minus $12$).}

Iškeliame bendrus dauginamuosius iš kiekvieno skliausto (iš pirmo $x^2$, iš antro $4$):
\[ x^2(x - 3) - 4(x - 3) \]
Dabar matome, kad abu masyvūs dėmenys turi bendrą „bloką“ $(x - 3)$. Iškeliame jį:
\[ (x - 3)(x^2 - 4) \]
Ar tai pabaiga? Olimpiadose visada patikrinkite, ar skliaustuose neliko formulių. Antrasis skliaustas yra akivaizdus kvadratų skirtumas ($x^2 - 2^2$)! Išskaidome jį iki galo:
\[ (x - 3)(x - 2)(x + 2) \hfill \square \]

\textbf{b)} Šiame reiškinyje vėl 4 nariai. Pirmieji du nariai iškart „prašosi“ būti išskaidyti pagal formulę:
\[ (a^2 - b^2) + (5a - 5b) \]
Išskaidome pirmąjį skliaustą, o iš antrojo iškeliame 5:
\[ (a - b)(a + b) + 5(a - b) \]
Matome, kad susidarė du dideli dėmenys, turintys bendrą skliaustą $(a - b)$. Iškeliame jį:
\[ (a - b)(a + b + 5) \hfill \square \]


\begin{example}[Kvadratų skirtumo grandinė]
Apskaičiuokite: $100^2 - 99^2 + 98^2 - 97^2 + 96^2 - 95^2 + \dots + 2^2 - 1^2$.
\end{example}

\textbf{Sprendimas:}\\
Niekas olimpiadose neprašo jūsų skaičiuoti $10000 - 9801$ ranka. Grupuojame po du:
\[ (100^2 - 99^2) + (98^2 - 97^2) + \dots + (2^2 - 1^2) \]
Išskaidome kiekvieną porą pagal formulę $a^2 - b^2 = (a-b)(a+b)$:
\[ (100-99)(100+99) + (98-97)(98+97) + \dots + (2-1)(2+1) \]
Pastebėkite, kad kiekviename pirmajame skliauste (skirtume) gauname lygiai $1$! 
Vadinasi, reiškinys virsta į:
\[ 1 \cdot (100+99) \quad + \quad 1 \cdot (98+97) \quad + \quad \dots \quad + \quad 1 \cdot (2+1) \]
Kadangi daugyba iš $1$ nieko nekeičia, mes galime tiesiog nutirpdyti skliaustus:
\[ 100 + 99 + 98 + 97 + \dots + 2 + 1 \]
Tai tiesiog visų natūraliųjų skaičių nuo 1 iki 100 suma. Pritaikome Gauso sumos taisyklę: $\frac{100 \cdot 101}{2} = 5050$. \hfill $\square$


\begin{example}[Milžiniškų skaičių kubai]
Apskaičiuokite (be skaičiuotuvo): $\frac{2026^3 + 1}{2026^2 - 2025}$.
\end{example}
\textbf{Sprendimas:}\\
Matome kubus, nes $1 = 1^3$. Užuot skaičiavę milijardus, įsiveskime kintamąjį. 
Pažymėkime $2026 = x$. Tuomet vardiklyje esantis $2025$ yra $x - 1$.
Perrašykime reiškinį naudodami $x$:
\[ \frac{x^3 + 1^3}{x^2 - (x - 1)} \quad = \quad \frac{x^3 + 1}{x^2 - x + 1} \]
Dabar skaitikliui pritaikome kubų sumos formulę $a^3 + b^3 = (a+b)(a^2 - ab + b^2)$:
\[ x^3 + 1^3 = (x + 1)(x^2 - x \cdot 1 + 1^2) = (x + 1)(x^2 - x + 1) \]
Įstatome šį skaidinį atgal į trupmeną:
\[ \frac{(x + 1)\mathbf{(x^2 - x + 1)}}{\mathbf{x^2 - x + 1}} \]
Masyvus antrasis skliaustas idealiai susiprastina su vardikliu! Lieka tik:
\[ x + 1 \]
Kadangi $x = 2026$, atsakymas yra $2026 + 1 = \mathbf{2027}$. \hfill $\square$


\begin{example}[Dalumo įrodymas su Sophie Germain]
Įrodykite, kad skaičius $2026^4 + 4^{2026}$ yra sudėtinis (t.y. turi daugiau daliklių nei 1 ir jis pats).
\end{example}

\textbf{Sprendimas:}\\
Tai yra išraiška $a^4 + 4b^4$. Pabandykime atrasti tinkamus $a$ ir $b$.
Akivaizdu, kad $a = 2026$. 
O kaip su $4^{2026}$? Žinome, kad $4^{2026} = (2^2)^{2026} = 2^{4052}$. 
Išskirkime daugiklį $4$: $2^{4052} = 2^2 \cdot 2^{4050} = 4 \cdot (2^{1012})^4$.
Taigi pradinė išraiška yra:
\[ 2026^4 + 4 \cdot (2^{1012})^4 \]
Pagal Sophie Germain tapatybę, šį skaičių (kur $a = 2026$ ir $b = 2^{1012}$) galime išskaidyti į dviejų skliaustų sandaugą:
\[ (a^2 - 2ab + 2b^2)(a^2 + 2ab + 2b^2) \]
Kadangi nei $a$, nei $b$ nėra lygūs $0$, abu šie skliaustai bus natūralieji skaičiai, griežtai didesni už $1$. Jei skaičių galima užrašyti kaip dviejų skaičių (didesnių už 1) sandaugą, vadinasi, jis yra sudėtinis. \hfill $\square$

\vspace{0.5cm}

% ==========================================
% UŽDAVINIAI
% ==========================================
\newpage
\section{Uždaviniai savarankiškam sprendimui (18 iššūkių)}

\begin{enumerate}
    \renewcommand{\labelenumi}{\textbf{\theenumi.}}
    
    \item Apskaičiuokite (be skaičiuotuvo!): $\frac{2026^2 - 2025^2}{2026 + 2025}$.
    \item Apskaičiuokite: $2026^2 - 2027 \cdot 2025$.
    \item Išskaidykite dauginamaisiais: $x^3 - x^2 - 4x + 4$.
    \item Išskaidykite dauginamaisiais: $a^2 - b^2 - a - b$.
    
    \item Apskaičiuokite: $1^2 - 2^2 + 3^2 - 4^2 + 5^2 - \dots - 100^2$.
    \item Apskaičiuokite (be skaičiuotuvo): $\frac{77^3 + 23^3}{100} - 77 \cdot 23$.
    \item Suprastinkite trupmeną: $\frac{a^3 - 8}{a^2 + 2a + 4}$.
    \item Išskaidykite dauginamaisiais reiškinius:
    \begin{itemize}
        \item[a)] $x^4 + 4$ 
        \item[b)] $y^4 + y^2 + 25$
    \end{itemize}
    \item Reiškinį $x^6 - y^6$ išskaidykite dauginamaisiais.
    
    \item Duota, kad $x - y = 5$ ir $xy = 14$. Nespręsdami pačios lygties (neieškodami $x$ ir $y$), raskite reiškinio $x^2 + y^2$ reikšmę.
    \item Žinoma, kad $a + \frac{1}{a} = 3$. Raskite:
    \begin{itemize}
        \item[a)] $a^2 + \frac{1}{a^2}$
        \item[b)] $a^3 + \frac{1}{a^3}$
    \end{itemize}
    \item Išskaidykite reiškinį $a^3 + b^3 + a + b$. Ar visada, jei $a$ ir $b$ yra nelyginiai skaičiai, ši išraiška dalijasi iš $4$?
    
    \item Išskaidykite dauginamaisiais reiškinį $x^4 + 64$.
    \item Įrodykite, kad skaičius $3^{44} + 4^{29}$ yra sudėtinis.
    \item Ar egzistuoja toks pirminis skaičius $p$, kuriam esant reiškinys $p^4 + 4$ taip pat yra pirminis skaičius?
    \item Išskirdami pilną kvadratą įrodykite, kad reiškinys $x^2 - 6x + 13$ niekada negali būti neigiamas. Kokia yra pati mažiausia galima šio reiškinio reikšmė ir su kokiu $x$ ji pasiekiama?
    \item Raskite visus natūraliuosius skaičius $n$, su kuriais reiškinys $n^4 + 4^n$ yra pirminis skaičius.
    \item Įrodykite, kad bet kokiam natūraliajam $n > 1$, skaičius $n^4 + 4n^3 + 11n^2 + 14n + 8$ yra sudėtinis.
\end{enumerate}

% ==========================================
% ATSAKYMAI
% ==========================================
\newpage
\section*{Atsakymai}

\begin{enumerate}
    \item $1$
    \item $1$
    \item $(x-1)(x-2)(x+2)$
    \item $(a+b)(a-b-1)$
    \item $-5050$
    \item $2916$
    \item $a-2$
    \item a) $(x^2-2x+2)(x^2+2x+2)$; b) $(y^2-3y+5)(y^2+3y+5)$
    \item $(x-y)(x+y)(x^2-xy+y^2)(x^2+xy+y^2)$
    \item $53$
    \item a) $7$; b) $18$
    \item Išskaidymas: $(a+b)(a^2-ab+b^2+1)$. Taip, visada (nes gaunama dviejų lyginių skaičių sandauga).
    \item $(x^2-4x+8)(x^2+4x+8)$
    \item Įrodyta (remiantis Sophie Germain tapatybe, abu daugikliai didesni už 1).
    \item Ne. (Vienintelis kandidatas $p=1$ nėra pirminis skaičius).
    \item Mažiausia galima reikšmė yra $4$, ji pasiekiama, kai $x=3$.
    \item $n=1$
    \item Įrodyta (su visais $n$ reiškinio reikšmė yra lyginis skaičius, o kai $n>1$, tas skaičius griežtai didesnis už $2$, todėl jis sudėtinis).
\end{enumerate}

\end{document}